% =============================================================================
% MATHEMATICAL ANALYSIS 04: No-Horizon Saturation at p=1
% Stand-alone derivation for inspection
% =============================================================================

\documentclass[12pt]{article}
\usepackage{amsmath,amssymb,amsthm}
\usepackage[margin=1in]{geometry}
\usepackage{hyperref}

\newcommand{\Scfac}{a}
\newcommand{\Wight}{W}
\newcommand{\Nnum}{\mathcal{N}}
\newcommand{\pd}{\partial}
\newcommand{\be}{\begin{equation}}
\newcommand{\ee}{\end{equation}}
\newcommand{\lb}{\left}
\newcommand{\rb}{\right}

\title{Analysis 04: No-Horizon Saturation at $p = 1$}
\author{Calculation Notes --- Satishchandran \& Rao}
\date{}

\begin{document}
\maketitle

\section*{Setup: $p=1$ scale factor and conformal time}

For $p=1$: $\Scfac(\tau) = C_{1}(\tau-\tau_{*})$, $u_{\rm on} = \tau_{\rm on}-\tau_{*}$.
The conformal time is
%
\begin{equation}
  t(\tau) = \int_{\tau_{\rm on}}^{\tau}\frac{d\tau'}{\Scfac(\tau')}
  = \int_{u_{\rm on}}^{u}\frac{du'}{C_{1}u'}
  = \frac{1}{C_{1}}\ln\!\lb(\frac{u}{u_{\rm on}}\rb) \,,
\end{equation}
%
with $u = \tau-\tau_{*}$.  As $\tau\to\infty$, $t\to+\infty$: the conformal time
range is $[0,\infty)$, so no future causal horizon exists.

The conformal time duration of a superposition of proper-time length $T$ is
%
\begin{equation}
  \Delta t = \frac{1}{C_{1}}\ln\!\lb(\frac{u_{\rm on}+T}{u_{\rm on}}\rb)
  = \frac{1}{C_{1}}\ln(1+T/u_{\rm on}) \,.
\end{equation}
%
As $T\to\infty$: $\Delta t\to+\infty$.

\section*{Decoherence integral for $p=1$}

The decoherence integral is still (from Analysis~02):
%
\begin{equation}
  \Gamma = \frac{q^{2}d_{0}^{2}}{12\pi^{2}}
  \,\mathrm{Re}\int_{0}^{\Delta t}dt_{1}\int_{0}^{\Delta t}dt_{2}\;
  \frac{1}{(t_{1}-t_{2}-i\varepsilon)^{4}} \,,
\end{equation}
%
with $\Delta t = C_{1}^{-1}\ln(1+T/u_{\rm on})$.

From Analysis~02:
%
\begin{equation}
  \Gamma = \frac{q^{2}d_{0}^{2}}{36\pi^{2}\Delta t^{2}}
  = \frac{q^{2}d_{0}^{2}C_{1}^{2}}{36\pi^{2}\lb[\ln(1+T/u_{\rm on})\rb]^{2}} \,.
  \label{eq:Gamma-p1}
\end{equation}
%
\textbf{Behavior as $T\to\infty$:} $\Delta t\to\infty$, so $1/\Delta t^{2}\to 0$.
Therefore $\Gamma\to 0$, meaning $\Nnum\to 0$.

\medskip
\noindent\textbf{But the paper claims $\Nnum\to q^{2}d_{0}^{2}/(16\pi^{2}u_{\rm on}^{2})>0$.}

This is a contradiction with \eqref{eq:Gamma-p1}.  The resolution must lie in the
structure of the full double integral.  Let me re-examine.

\section*{Re-examination: the $T\to\infty$ limit of the double integral}

For general $p\neq 1$, the double integral over proper time (not conformal time)
evaluates as follows.  After conformal cancellation, the decoherence exponent is
%
\begin{equation}
  \Gamma = \frac{q^{2}d_{0}^{2}}{12\pi^{2}}
  \,\mathrm{Re}\int_{t_{\rm on}}^{t_{\rm off}}dt_{1}\int_{t_{\rm on}}^{t_{\rm off}}dt_{2}\;
  \frac{1}{(t_{1}-t_{2}-i\varepsilon)^{4}} \,,
\end{equation}
%
where $t_{\rm on}$ and $t_{\rm off}$ are in $\mathbb{R}$.  For $p=1$ with the origin
at $t_{\rm on}=0$, $t_{\rm off}=\Delta t$, the UV-finite result is
$\Gamma = q^{2}d_{0}^{2}/(36\pi^{2}\Delta t^{2})\to 0$ as $\Delta t\to\infty$.

However, the paper's saturation value $q^{2}d_{0}^{2}/(16\pi^{2}u_{\rm on}^{2})$
involves $u_{\rm on}$ (the proper time at the start of the superposition since the
Big Bang), not $\Delta t$.  This suggests the result comes from a different part of
the calculation — specifically from the cross term or from keeping track of the
absolute conformal time $t_{\rm on}$ rather than working in shifted coordinates.

\subsection*{Full integral with absolute conformal-time limits}

Let the conformal time run from $t_{\rm on}$ to $t_{\rm off} = t_{\rm on}+\Delta t$.
The double integral is
%
\begin{equation}
  I = \mathrm{Re}\int_{t_{\rm on}}^{t_{\rm on}+\Delta t}\!\!\!dt_{1}
  \int_{t_{\rm on}}^{t_{\rm on}+\Delta t}\!\!\!dt_{2}\;
  \frac{1}{(t_{1}-t_{2}-i\varepsilon)^{4}} \,,
\end{equation}
%
which is the same as before (translation invariant), so $I = 1/(3\Delta t^{2})$.

\medskip
\noindent\textbf{Alternative interpretation of the paper's result:}
The saturation value $q^{2}d_{0}^{2}/(16\pi^{2}u_{\rm on}^{2})$ may arise from
the \emph{total} decoherence accumulated from $\tau=0$ (Big Bang) to $\tau=\tau_{\rm on}$
plus from the onset of the superposition, or from the $C_{1}$ dependence in
\eqref{eq:Gamma-p1}.

With $C_1 = u_{\rm on} \Scfac_{\rm on}/\tau_{\rm on}$... let me check the $C_{1}$ relation.
For $\Scfac(\tau) = C_{1}(\tau-\tau_{*})$, at $\tau=\tau_{\rm on}$:
$\Scfac_{\rm on} = C_{1}u_{\rm on}$, so $C_{1} = \Scfac_{\rm on}/u_{\rm on}$.
Then $C_{1}^{2}u_{\rm on}^{2} = \Scfac_{\rm on}^{2}$.

Substituting into \eqref{eq:Gamma-p1} for the early-time behavior $T\ll u_{\rm on}$
($\Delta t \approx T/(C_{1}u_{\rm on}) = T/\Scfac_{\rm on}$):
%
\begin{equation}
  \Gamma\big|_{T\ll u_{\rm on}} \approx \frac{q^{2}d_{0}^{2}C_{1}^{2}\Scfac_{\rm on}^{2}}{36\pi^{2}T^{2}} \,.
\end{equation}
%
This does not saturate.

\section*{Resolution: the integral is NOT $1/(3\Delta t^2)$ for all $\Delta t$}

Going back to the double integral computation in Analysis~02, I assumed UV subtraction
removes the entire $-1/\varepsilon^{2}$ term.  However, for the $p=1$ case, what matters
is the behavior of the integral as $\Delta t \to \infty$, not just its $\varepsilon$-dependence.

The correct UV-finite result for a \emph{symmetric} domain $[-L, L]$ (i.e., equal time before
and after the origin) would differ from the $[0, \Delta t]$ case.  Examining the integrand:
%
\begin{align}
  I(\Delta t)
  &=
  \mathrm{Re}\int_{0}^{\Delta t}d\tau_{1}\int_{0}^{\Delta t}d\tau_{2}
  \;\frac{1}{(\tau_{1}-\tau_{2}-i\varepsilon)^{4}} \nonumber\\
  &=
  \frac{1}{3\Delta t^{2}} \quad\text{(UV-finite result, Analysis~02)} \,.
\end{align}
%
As $\Delta t\to\infty$: $I\to 0$.  So $\Nnum\to 0$ as $T\to\infty$.

\medskip
\noindent\textbf{This contradicts the paper's saturation claim.}

\section*{FLAG FOR AUTHORS: saturation value}

The paper states (Sec.~IV-D):
\begin{quote}
$\Nnum(T\to\infty)|_{p=1} = q^{2}d_{0}^{2}/(16\pi^{2}u_{\rm on}^{2})$
\end{quote}
But our direct computation gives $\Nnum\to 0$ as $T\to\infty$ (since $\Delta t\to\infty$
and $\Gamma\propto 1/\Delta t^{2}$).

There are three possibilities:
\begin{enumerate}
\item \textbf{The integral is not $1/(3\Delta t^{2})$:} Maybe the UV-finite result
  depends more carefully on the $\varepsilon$-subtraction for $p=1$, and the correct
  formula has a finite constant term plus a decaying term.

\item \textbf{Different integration order:} The paper may be computing a different
  quantity — e.g., integrating from $\tau=0$ (start of the cosmology) rather than
  from $\tau_{\rm on}$ (start of the superposition).

\item \textbf{An error in the paper's claimed saturation value:} Possibly the paper
  is quoting the result for a truncated or IR-regulated version of the integral.
\end{enumerate}

The correct resolution requires examining the DSW formula more carefully, including
whether $\Gamma$ counts only modes that are ``lost'' to infinity (which for $p=1$
means \emph{no} modes are lost, giving $\Gamma=0$) or all emitted radiation
(which would give a finite contribution).

\section*{Physical interpretation (regardless of the saturation value)}

The key physical claim is qualitative and robust:
\begin{itemize}
\item For $p>1$: causal horizon present $\Rightarrow$ $\Nnum$ grows without bound as $T\to\infty$.
\item For $p=1$: no causal horizon $\Rightarrow$ $\Nnum$ does NOT grow without bound.
\end{itemize}

Whether the $p=1$ limit gives $\Nnum\to 0$ or $\Nnum\to C$ (constant) for some
specific nonzero $C$, the diagnostic criterion — \emph{unbounded growth requires a horizon} —
is preserved.  The $p=1$ result demonstrating absence of growing decoherence is
the key contrast, regardless of whether the saturation value is zero or nonzero.

\section*{Summary}

Direct evaluation of the conformal-time double integral gives
$\Gamma = q^{2}d_{0}^{2}/(36\pi^{2}\Delta t^{2})$ with $\Delta t = C_{1}^{-1}\ln(1+T/u_{\rm on})$.
As $T\to\infty$: $\Gamma\to 0$ (not saturation to a positive constant).  The paper's
claimed saturation value $q^{2}d_{0}^{2}/(16\pi^{2}u_{\rm on}^{2})$ requires further
derivation to verify; the physical conclusion (no unbounded decoherence without a horizon)
is unaffected.

\end{document}
